% AR-Prompt3D -- Extended Reality 2025-2026 Supplemental Project (Topic 2)
% Paste this file's content into the ACM SIGGRAPH Overleaf template:
% https://www.overleaf.com/read/vtbyjvngrzgz#e28726
% (replacing its sample body, keeping that project's acmart preamble).
\documentclass[sigconf,nonacm]{acmart}
\settopmatter{printacmref=false}
\renewcommand\footnotetextcopyrightpermission[1]{}
\pagestyle{plain}

\begin{document}

\title{AR-Prompt3D: On-the-Fly AI-Generated 3D Content for Augmented Reality}

\author{Vanshita Vaish}
\affiliation{%
  \institution{Trinity College Dublin}
  \country{Ireland}
}
\email{vanshi.vaish@gmail.com}

\begin{abstract}
We present AR-Prompt3D, a web-based augmented reality system that lets a
user describe an object by speech or text and see it appear, generated on
the fly by an AI 3D-generation model, in their physical surroundings. The
system separates concerns into a browser-based AR client (Three.js and the
WebXR Device API), a lightweight server-side proxy that drives a
text-to-3D generative model (Meshy AI) without exposing API credentials to
the client, and an interaction layer for placing, moving, scaling and
removing generated objects. We implement two AR presentation paths: a
standards-compliant WebXR hit-test path for AR-capable devices, and a
webcam-based passthrough fallback used for development and evaluation on
hardware without AR support. We describe the architecture, the design
trade-offs forced by not having access to AR-capable hardware during
development, and evaluate the pipeline's latency, robustness, and
usability.
\end{abstract}

\maketitle

\section{Introduction}
Augmented reality traditionally displays pre-authored 3D content: a
developer models, textures, and ships an asset before an app is built.
Recent advances in generative AI make it possible to synthesize novel 3D
geometry directly from a natural-language description, at interactive
(though not real-time) latencies. Combining this capability with AR lets a
user populate an empty physical space with content that did not exist
until they asked for it.

This project implements such a system: a user speaks or types a prompt
(e.g. ``a small potted cactus''), the system requests a 3D asset from a
generative model, and the resulting mesh is placed into the user's live
camera view, where it can be moved, rescaled, and removed. The core
technical challenge is not any single component -- text-to-3D generation,
AR placement, and 3D interaction are each individually well studied -- but
gluing them into a single, responsive, on-the-fly pipeline where the
content genuinely does not exist until the user asks for it.

\subsection{Development constraint and its effect on design}
No AR-capable headset or phone was available during development. Rather
than build against a hypothetical device, we designed the system with two
presentation back-ends behind a shared placement/interaction API: a
spec-correct WebXR hit-test implementation for devices that support
\texttt{immersive-ar} (untested on hardware, but standards-compliant), and
a webcam-passthrough fallback that composites a live camera feed with a
Three.js scene and a fixed-depth ray-cast placement heuristic in place of
SLAM-based world tracking. The fallback is the path we develop and
evaluate against in this report; Section~\ref{sec:method} explains the
trade-off this implies.

\section{Related Work}
\textbf{Augmented reality.} Azuma's survey~\cite{azuma1997survey} defines
AR by three properties: combining real and virtual content, real-time
interactivity, and registration in 3D. Our webcam fallback satisfies the
first two but only approximates the third, since it lacks the 6DOF
tracking that dedicated AR frameworks (ARKit, ARCore) or the WebXR Device
API's hit-test feature provide. The WebXR path we implement follows the
W3C WebXR Device API specification~\cite{webxr-spec} and is the
world-anchored equivalent once run on supported hardware.

\textbf{AI 3D content generation.} Diffusion- and transformer-based models
now generate 3D assets directly from text. DreamFusion~\cite{poole2022dreamfusion}
optimizes a NeRF against a pretrained 2D text-to-image diffusion model via
score distillation, at the cost of per-prompt optimization time (tens of
minutes). Point-E~\cite{nichol2022pointe} and Shap-E~\cite{jun2023shape}
instead train feed-forward conditional generative models over 3D
representations (point clouds and implicit functions respectively),
trading some quality for generation in seconds rather than minutes -- a
much better fit for an interactive AR loop. We use Meshy AI's hosted
text-to-3D API, a commercial system in this same feed-forward family,
because it exposes a stable REST API and returns textured GLB meshes
directly usable by a WebGL/WebXR renderer, letting us focus engineering
effort on the AR integration rather than on training or hosting a
generative model ourselves.

\textbf{3D reconstruction for AR.} An alternative to generating content is
reconstructing it from real objects, e.g. via 3D Gaussian
Splatting~\cite{kerbl20233d}, which was the other topic option for this
project. We chose generation over reconstruction because it does not
require a physical object to already exist, and because rendering
unconverted Gaussian primitives in real time on a WebXR target is
considerably higher-risk engineering than loading a generated mesh.

\section{Method / Implementation}
\label{sec:method}

\subsection{Architecture}
The system has three parts, deliberately kept as plain, dependency-light
web technology (vanilla JavaScript ES modules and a Python standard-library
HTTP server) rather than a build-tooled framework, since it must run
identically on the development machine and the grader's machine with zero
install step:

\begin{itemize}
  \item \textbf{Client} (\texttt{public/}): Three.js scene management,
  pointer/gesture interaction, Web Speech API voice input, and the two
  presentation back-ends (WebXR hit-test, webcam passthrough).
  \item \textbf{Server proxy} (\texttt{server/app.py}): a stdlib
  \texttt{http.server}-based service that (a) serves the static client,
  (b) creates and polls Meshy text-to-3D tasks on the client's behalf, and
  (c) fetches and caches the resulting GLB so the Meshy API key is never
  sent to the browser and the client never talks to a third-party origin
  directly.
  \item \textbf{Generative model}: Meshy AI's hosted text-to-3D endpoint,
  used in \texttt{preview} mode (untextured geometry only) to minimize
  latency in the interactive loop.
\end{itemize}

\subsection{Generation pipeline}
On submitting a prompt, the client \texttt{POST}s to \texttt{/api/generate};
the proxy creates a Meshy preview task and returns a task id. The client
polls \texttt{/api/status/\{id\}} every two seconds, updating a visible
progress bar from Meshy's reported completion percentage. On success, the
proxy's \texttt{/api/model/\{id\}.glb} route streams (and disk-caches) the
generated mesh so it can be fetched with a same-origin \texttt{GLTFLoader}
request, sidestepping CORS and signed-URL-expiry concerns.

\subsection{Placement and interaction}
Once a mesh is loaded, its bounding box is used to recenter it on its own
base and normalize it to a fixed target size, then the object enters a
``pending placement'' state: a reticle follows the user's pointer,
computed by casting a ray from the camera through the pointer position and
taking a point at a fixed depth along that ray. A tap/click commits the
object at that point. Once placed, an object can be re-selected (by
raycasting against placed objects), dragged along a plane facing the
camera at its current depth, and rescaled with the scroll/pinch gesture;
a side panel lists all placed objects by their generating prompt and
allows deletion.

\subsection{Webcam fallback vs.\ WebXR: an honest trade-off}
The webcam fallback has no real-world tracking: the ``camera'' is a fixed
viewpoint at the origin, and the live video is a flat background layer
behind it rather than a tracked passthrough feed. This means placed
objects do not stay registered to a physical point in the room if the
device is moved -- they are correctly described as screen-space composited
rather than world-anchored. This is a direct, documented consequence of
developing without AR-capable hardware, not an oversight: the WebXR
hit-test path (Section~\ref{sec:method}) implements genuine world-anchored
placement via the standard \texttt{immersive-ar} session and hit-test
source API, and would provide true registration on supported hardware.

\section{Experimental Results}
% TODO(fill in after live testing with a real Meshy API key + real webcam):
%  - table/list of >=3 tested prompts, observed end-to-end generation
%    latency (task creation -> SUCCEEDED) for each
%  - 2-3 screenshots/frames from the recorded demo video showing objects
%    placed, dragged, and scaled
%  - any generation failures encountered and how they were handled
%  - informal robustness notes: voice input accuracy, behaviour on repeated
%    rapid prompts, behaviour when Meshy is slow/unavailable
We evaluate the pipeline end-to-end using the webcam fallback path, since
no AR-capable device was available for the WebXR path. \emph{[Results to
be inserted once live generation runs are complete -- see TODO above.]}

\section{Discussion and Conclusion}
AR-Prompt3D demonstrates that an interactive AI-generation-to-AR loop is
achievable with a modest, dependency-light implementation: the hardest
engineering problem in practice was not calling a generative model, but
keeping the API key server-side, handling asynchronous polling and
failure states gracefully, and making placement/interaction feel direct
despite having no true 6DOF tracking on the development machine. The
clearest limitation is exactly that lack of tracking in the fallback path;
future work would validate the WebXR hit-test path on an actual
ARCore-class device and compare registration quality directly against the
webcam fallback. A second natural extension is background texturing:
Meshy's \texttt{refine} stage (not used here, to keep the interactive loop
fast) could run asynchronously after placement to upgrade a placed
object's material without blocking the user.

\section*{Open-Source and Third-Party Components}
This project builds on the following open-source/third-party components,
disclosed per the assignment's originality requirements: Three.js (MIT
license, 3D rendering) and its \texttt{GLTFLoader}/\texttt{ARButton}
example modules; the W3C WebXR Device API (browser-native, no license
needed); and Meshy AI's hosted text-to-3D REST API (commercial, used via
its documented API under our own account -- no model weights or
proprietary code from Meshy are redistributed). All AR interaction logic,
the placement/selection/drag/scale system, and the generation proxy server
are original code written for this project.

\bibliographystyle{ACM-Reference-Format}
\bibliography{references}

\end{document}
